\subsubsection{CGNS}
\label{SEC:Files:Grid_Input:CGNS}

CGNS is the recommended format to use for grid files with GFR.
%
The primary reason for this is because CGNS is one of the few commonly used CFD
grid formats that supports high-order curved elements.
%
The Gmsh file format also supports high-order curved elements, but the GFR
Gmsh grid reader is severely limited and out-of-date.
%
The CGNS file format is very flexible in allowing for multiple ways to store the
same computational mesh, which can make it difficult to create a universal CGNS grid
reader that never fails.
%
All of the CGNS grids used thus far with GFR have been generated by the Cadence
Fidelity Pointwise software.
%
Therefore, grids generated by Pointwise should have the highest success at
being read and successfully processed by GFR.
%
CGNS files generated by other grid generation software have not been tested
for compatibility with GFR.

%%%%%%%%%%%%%%%%%%%%%%%%%%%%%%%%%%%%%%%%%%%%%%%%%%%%%%%%%%%%%%%%%%%%%%%%%%%%%%%%

When the solid body surfaces in a flow contain nonlinear curvature, using a
tessellation of linear segments to approximate this curvature --- as is done
with linear grids used for traditional 2\supnd-order finite-volume methods ---
results in a misrepresentation of the surface for high-order finite-element-type
methods, \forexample FR\slash DG.
%
Without any curvature information, these high-order methods will treat the
linear tessellation of the curved surface as a series of flat plates, with each
new linear segment of the surface resulting in a new compression or expansion
corner.
%
This linear approximation is sufficient for 2\supnd-order finite-volume methods
because 1) the compression\slash expansion angle between linear segments is
generally smaller for these methods since they generally require
higher-resolution grids compared to higher-order methods, and 2) the error from
the numerical discretization generally provides enough numerical dissipation to
damp out any non-physical waves that might emanate from the intersection of two
linear segments.
%
All of this is to say that high-order curved meshes are critical for high-order
methods to achieve their accuracy potential for flows involving solid body
surfaces with nonlinear curvature.

%%%%%%%%%%%%%%%%%%%%%%%%%%%%%%%%%%%%%%%%%%%%%%%%%%%%%%%%%%%%%%%%%%%%%%%%%%%%%%%%

We have used Pointwise to generate grids for GFR because it is one of the few
available grid generation tools that supports unstructured,
high-order curved meshes.
%
Pointwise provides both written documentation and online videos demonstrating
how to use the software to generate high-order meshes, so that information will
not be covered here.
%
The one important thing is to make sure any curved boundary surfaces are
accurately represented as a Database surface in Pointwise, and that the Domain
meshes on these surfaces have been projected to their respective Databases.
%
The remainder of this section will cover the steps in the Pointwise grid
generation process that affects the use of the grid within GFR.

%%%%%%%%%%%%%%%%%%%%%%%%%%%%%%%%%%%%%%%%%%%%%%%%%%%%%%%%%%%%%%%%%%%%%%%%%%%%%%%%

\subsubsubsection{Setting Boundary Conditions in Pointwise}
\label{SEC:Files:Grid_Input:CGNS:PWBC}

The primary task required to effectively use a CGNS grid exported from Pointwise
with GFR is setting the boundary conditions for the grid.
%
Setting the boundary conditions in Pointwise starts with accessing the
\texttt{"Set BC"} dialog box by selecting \texttt{"Set Boundary
Conditions\ldots"} under the \texttt{"CAE"} menu.
%
Figures \ref{FIG:Files:PW:Spacefan:BC} and \ref{FIG:Files:PW:THX3:BC} show
examples of the \texttt{"Set BC"} dialog box for two different cases previously
run with GFR.
%
Images of the respective grids are also shown alongside the set of
boundary conditions in figures \ref{FIG:Files:PW:Spacefan:Geometry} and
\ref{FIG:Files:PW:THX3:Geometry} as a reference.

%%%%%%%%%%%%%%%%%%%%%%%%%%%%%%%%%%%%%%%%%%%%%%%%%%%%%%%%%%%%%%%%%%%%%%%%%%%%%%%%

In the \texttt{"Set BC"} dialog box, the third column titled \texttt{"Name"} and
the fourth column titled \texttt{"CAE Type"} are the user inputs that are used
to configure the boundary conditions for use with GFR.
%
The \texttt{"Name"} column allows you to input a string to give each boundary
condition a name.
%
The string provided in Pointwise as the name for a boundary condition is used to match
against one of the boundary conditions in the GFR input file, specifically
against the inputs \BCinpBC[nmlBC::BCInput/Name][i]{Name}.
%
The \texttt{"CAE Type"} column provides a drop down list of different boundary
condition types supported by the CGNS standard.
%
Note that certain boundary condition types supported by the CGNS standard are
not supported by GFR, but GFR also supports boundary condition types not
supported by the CGNS standard.
%
The boundary condition options in the drop down list are given in the first
column of \tref{TBL:PW_BC_GFR_Match}.
%
Ideally, the GFR input file will also provide the respective
\BCinpBC[nmlBC::BCInput/Name][i]{Set\_BC\_Value} to specify the type of
boundary condition to use for the boundary.
%
If \inptext{Set\_BC\_Value} is not
specified, the option specified under \texttt{"CAE Type"} is used to set the
type of boundary condition to use for the boundary, and
\tref{TBL:PW_BC_GFR_Match} lists the GFR boundary condition assigned based on
the \texttt{"CAE Type"}.
%
%###############################################################################
%
%\begin{figure}[p]
\begin{figure}[b!]
 %\setlength{\mpwidth}{0.52\textwidth}
  \setlength{\imagewidth}{1.0\textwidth}
  \setlength{\imageheight}{0.33\textheight}
  \centering
  \subfig{\imagewidth}{
    \insertimage[height=\imageheight]{pointwise_bcs/spacefan_bcs}
    \caption{Pointwise \texttt{"Set BC"} dialog box}
    \label{FIG:Files:PW:Spacefan:BC}
  }
  \\
 %\vspace{10pt}
  \subfig{\imagewidth}{
    \insertimage[height=\imageheight]{pointwise_bcs/spacefan_geometry}
    \caption{Computational geometry}
    \label{FIG:Files:PW:Spacefan:Geometry}
  }
% \caption{Configuring the boundary conditions in Pointwise for the spacecraft
%          cabin ventilation fan grid. The string in the third column titled
%          \texttt{"Name"} is used to match against
%          \BCinptext{Name} in the GFR input file. The selection
%          in the fourth column titled \texttt{"CAE Type"} is used to set the
%          boundary condition value based on \tref{TBL:PW_BC_GFR_Match} if the
%          input \BCinptext{Set\_BC\_Value} is not specified for
%          a boundary condition with the respective name.}
  \caption[Pointwise boundary condition settings for the spacecraft
           cabin ventilation fan grid.]%
          {Pointwise boundary condition settings for the spacecraft
           cabin ventilation fan grid.
           The string in the third column titled \texttt{"Name"} is used to
           match against \BCinptext{Name} in the GFR input file.
           The selection in the fourth column titled \texttt{"CAE Type"} is
           used to set the boundary condition value based on
           \tref{TBL:PW_BC_GFR_Match} if the input \BCinptext{Set\_BC\_Value}
           is not specified for a boundary condition with the respective name.}
  \label{FIG:Files:PW:Spacefan}
\end{figure}

%%%%%%%%%%%%%%%%%%%%%%%%%%%%%%%%%%%%%%%%%%%%%%%%%%%%%%%%%%%%%%%%%%%%%%%%%%%%%%%%

\noindent\begin{minipage}{\textwidth}
\begin{nmldescitem}[%
   fontupper=\ttfamily\footnotesize%
  ,listing options={escapeinside={!@}{!@}}%
]
&nml_bc
  BC_Input(1)%Name           = "Inflow"
  BC_Input(1)%Set_BC_Value   = 2           ! Subsonic inflow
  BC_Input(1)%Total_Pressure = 102070.0    ! [Pa] (101992.2 static)
  BC_Input(1)%Temperature    = 294.3       ! [K]
  BC_Input(1)%Mach           = 0.033       ! Mach number !@\vspace{3.9pt}!@
  BC_Input(2)%Name           = "Outflow"
  BC_Input(2)%Set_BC_Value   = 12          ! Subsonic outflow
  BC_Input(2)%Pressure       = 102553.0    ! [Pa] (1.0055 * Inflow Pressure)
  BC_Input(2)%Temperature    = 294.3       ! [K]
  BC_Input(2)%Mach           = 0.033       ! Mach number !@\vspace{3.9pt}!@
  BC_Input(3)%Name                = "WallHub"
  BC_Input(3)%Set_BC_Value        = 52     ! Adiabatic wall
  BC_Input(3)%Is_Specified_In_MRF = T      ! Moving reference frame
  BC_Input(3)%Rotation_Rate(1)    = 0.0    ! [RPM] !@\vspace{3.9pt}!@
  BC_Input(4)%Name                = "WallHubSponge"
  BC_Input(4)%Set_BC_Value        = 52     ! Adiabatic wall
  BC_Input(4)%Is_Specified_In_MRF = T      ! Moving reference frame
  BC_Input(4)%Rotation_Rate(1)    = 0.0    ! [RPM] !@\vspace{3.9pt}!@
  BC_Input(5)%Name                = "WallCaseSponge"
  BC_Input(5)%Set_BC_Value        = 52     ! Adiabatic wall
  BC_Input(5)%Is_Specified_In_MRF = F      ! Fixed reference frame
  BC_Input(5)%Rotation_Rate(1)    = 0.0    ! [RPM] !@\vspace{3.9pt}!@
  BC_Input(6)%Name                = "WallCase"
  BC_Input(6)%Set_BC_Value        = 52     ! Adiabatic wall
  BC_Input(6)%Is_Specified_In_MRF = F      ! Fixed reference frame
  BC_Input(6)%Rotation_Rate(1)    = 0.0    ! [RPM] !@\vspace{3.9pt}!@
  BC_Input(7)%Name                = "WallBlade"
  BC_Input(7)%Set_BC_Value        = 52     ! Adiabatic wall
  BC_Input(7)%Is_Specified_In_MRF = T      ! Moving reference frame
  BC_Input(7)%Rotation_Rate(1)    = 0.0    ! [RPM] !@\vspace{3.9pt}!@
  BC_Input(8)%Name                      = "ThetaMin"
  BC_Input(8)%Set_BC_Value              = 81    ! Periodic BC
  BC_Input(8)%Periodic%Matching_BC_Name = "ThetaMax"
  BC_Input(8)%Periodic%Matching_BC_Idx  = 9
  BC_Input(8)%Periodic%RotationAngle(1) = 40.0  !  [deg] !@\vspace{3.9pt}!@
  BC_Input(9)%Name                      = "ThetaMax"
  BC_Input(9)%Set_BC_Value              = 81     ! Periodic BC
  BC_Input(9)%Periodic%Matching_BC_Name = "ThetaMin"
  BC_Input(9)%Periodic%Matching_BC_Idx  = 8
  BC_Input(9)%Periodic%RotationAngle(1) = -40.0  !  [deg] !@\vspace{3.9pt}!@
  BC_Input(10)%Name                      = "ThetaMinSponge"
  BC_Input(10)%Set_BC_Value              = 81     ! Periodic BC
  BC_Input(10)%Periodic%Matching_BC_Name = "ThetaMaxSponge"
  BC_Input(10)%Periodic%Matching_BC_Idx  = 11
  BC_Input(10)%Periodic%RotationAngle(1) = 40.0   ! [deg] !@\vspace{3.9pt}!@
  BC_Input(11)%Name                      = "ThetaMaxSponge"
  BC_Input(11)%Set_BC_Value              = 81     ! Periodic BC
  BC_Input(11)%Periodic%Matching_BC_Name = "ThetaMinSponge"
  BC_Input(11)%Periodic%Matching_BC_Idx  = 10
  BC_Input(11)%Periodic%RotationAngle(1) = -40.0  ! [deg]
/
\end{nmldescitem}
  \captionof{lstlisting}%
            [Example GFR inputs for the spacecraft cabin ventilation fan grid.]%
            {Example GFR inputs that correspond to the boundary condition
             configuration shown in \fref{FIG:Files:PW:Spacefan:BC} for the
             spacecraft cabin ventilation fan grid.}
  \label{TBL:Files:Spacefan:BCInputs}
\end{minipage}

%%%%%%%%%%%%%%%%%%%%%%%%%%%%%%%%%%%%%%%%%%%%%%%%%%%%%%%%%%%%%%%%%%%%%%%%%%%%%%%%

\begin{figure}[p]
  \setlength{\imagewidth}{1.0\textwidth}
  \setlength{\imageheight}{0.4\textheight}
  \centering
  \subfig{\imagewidth}{
    \insertimage[height=\imageheight]{pointwise_bcs/thx3_bcs}
    \caption{Pointwise \texttt{"Set BC"} dialog box}
    \label{FIG:Files:PW:THX3:BC}
  }
  \\
  \subfig{\imagewidth}{
    \insertimage[height=\imageheight]{pointwise_bcs/thx3_geometry}
    \caption{Computational geometry}
    \label{FIG:Files:PW:THX3:Geometry}
  }
% \caption{Configuring the boundary conditions in Pointwise for the THX3 grid.
%          The string in the third column titled \texttt{"Name"} is used to
%          match against \BCinptext[i]{Name} in the GFR input file.
%          The selection in the fourth column titled \texttt{"CAE Type"} is used
%          to set the boundary condition value based on
%          \tref{TBL:PW_BC_GFR_Match} if the input
%          \BCinptext[i]{Set\_BC\_Value} is not specified for a
%          boundary condition with the respective name.}
  \caption[Pointwise boundary condition settings for the THX3 grid.]%
          {Pointwise boundary condition settings for the THX3 grid.
           The string in the third column titled \texttt{"Name"} is used to
           match against \BCinptext{Name} in the GFR input file.
           The selection in the fourth column titled \texttt{"CAE Type"} is
           used to set the boundary condition value based on
           \tref{TBL:PW_BC_GFR_Match} if the input \BCinptext{Set\_BC\_Value}
           is not specified for a boundary condition with the respective name.}
  \label{FIG:Files:PW:THX3}
\end{figure}

%%%%%%%%%%%%%%%%%%%%%%%%%%%%%%%%%%%%%%%%%%%%%%%%%%%%%%%%%%%%%%%%%%%%%%%%%%%%%%%%

\noindent\begin{minipage}{\textwidth}
\begin{nmldescitem}[%
   fontupper=\ttfamily\footnotesize%
  ,listing options={escapeinside={!@}{!@}}%
]
&nml_bc
  BC_Input(1)%Name                    = "InflowNozzle"
  BC_Input(1)%Set_BC_Value            = 5            ! SEM inflow
  BC_Input(1)%Total_Pressure          = 108869.053   ! [Pa] (15.79 psi)
  BC_Input(1)%Total_Temperature       = 538.97       ! [K]  (970.146 R)
  BC_Input(1)%Mach                    = 0.3732       ! Inflow Mach number
  BC_Input(1)%SEM%Convective_Velocity = 84.723       ! [m/s] 1/2 of !@$U_\textsc{sem}$!@
  BC_Input(1)%SEM%Length_Scale        = 0.00030038   ! [m]   1/5 of !@$\delta_\textsc{bl}$!@
  BC_Input(1)%SEM%Profile_File        = "nozzle_exit.m-kg-s.match_P0-T0.dat" !@\vspace{5pt}!@
  BC_Input(2)%Name              = "InflowPlenum"
  BC_Input(2)%Set_BC_Value      = 2          ! Subsonic inflow
  BC_Input(2)%Total_Pressure    = 110500.0   ! [Pa] (16.03  psi)
  BC_Input(2)%Pressure          = 98595.03   ! [Pa] (14.3   psi)
  BC_Input(2)%Total_Temperature = 294.293    ! [K]  (529.7274 R)
  BC_Input(2)%AoA_Alpha         = 30.0       ! Inflow angle of attack !@\vspace{5pt}!@
  BC_Input(3)%Name              = "InflowExterior"
  BC_Input(3)%Set_BC_Value      = 2          ! Subsonic Inflow
  BC_Input(3)%Total_Pressure    = 98767.68   ! [Pa] (14.325 psi)
  BC_Input(3)%Total_Temperature = 294.591222 ! [K]  (530.2642 R)
  BC_Input(3)%Mach              = 0.05       ! Inflow Mach number !@\vspace{5pt}!@
  BC_Input(4)%Name              = "Farfield"
  BC_Input(4)%Set_BC_Value      = 21         ! Characteristic in/outflow
  BC_Input(4)%Total_Pressure    = 98767.68   ! [Pa] (14.325 psi)
  BC_Input(4)%Total_Temperature = 294.591222 ! [K]  (530.2642 R)
  BC_Input(4)%Mach              = 0.05       ! Inflow Mach number !@\vspace{5pt}!@
  BC_Input(5)%Name         = "Outflow"
  BC_Input(5)%Set_BC_Value = 12         ! Subsonic outflow
  BC_Input(5)%Pressure     = 98595.03   ! [Pa] (14.3 psi)
  BC_Input(5)%Temperature  = 294.444    ! [K]  (530.0 R)
  BC_Input(5)%VelocityX    = 169.4460   ! [m/s] !@$\delta_\textsc{bl}$!@ !@\vspace{5pt}!@
  BC_Input(6)%Name         = "WallPlate"
  BC_Input(6)%Set_BC_Value = 53           ! Isothermal wall
  BC_Input(6)%Wall_Temp    = 427.888889   ! [K] (770.2 R) !@\vspace{5pt}!@
  BC_Input(7)%Name         = "WallNozzleInterior"
  BC_Input(7)%Set_BC_Value = 52           ! Adiabatic wall !@\vspace{5pt}!@
  BC_Input(8)%Name         = "WallNozzleLip"
  BC_Input(8)%Set_BC_Value = 52           ! Adiabatic wall !@\vspace{5pt}!@
  BC_Input(9)%Name         = "WallNozzleExterior"
  BC_Input(9)%Set_BC_Value = 52           ! Adiabatic wall !@\vspace{5pt}!@
  BC_Input(10)%Name         = "WallPipe"
  BC_Input(10)%Set_BC_Value = 52          ! Adiabatic wall !@\vspace{5pt}!@
  BC_Input(11)%Name         = "WallPlenumSide"
  BC_Input(11)%Set_BC_Value = 52          ! Adiabatic wall !@\vspace{5pt}!@
  BC_Input(12)%Name         = "WallPlenumEnd"
  BC_Input(12)%Set_BC_Value = 52          ! Adiabatic wall !@\vspace{5pt}!@
  BC_Input(13)%Name         = "Bottom"
  BC_Input(13)%Set_BC_Value = 50          ! Slip wall
/
\end{nmldescitem}
  \captionof{lstlisting}%
            [Example GFR inputs for the THX3 grid.]%
            {Example GFR inputs that correspond to the boundary condition
             configuration shown in \fref{FIG:Files:PW:THX3:BC} for the THX3
             grid.}
  \label{TBL:Files:THX3:BCInputs}
\end{minipage}

%%%%%%%%%%%%%%%%%%%%%%%%%%%%%%%%%%%%%%%%%%%%%%%%%%%%%%%%%%%%%%%%%%%%%%%%%%%%%%%%

\begin{table}[tb]
  \centering
% \caption{List of CGNS boundary condition options and respective GFR boundary
%          condition assigned if \BCinptext[i]{Set\_BC\_Value} is
%          not used.}
  \caption[Assumed mapping between CGNS and GFR boundary conditions.]%
          {Assumed mapping of CGNS and GFR boundary conditions if
           \BCinptext[i]{Set\_BC\_Value} is not used. Unsupported CGNS boundary
           conditions will attempt to set the GFR boundary condition type based
           on the \quoterm{name match} method described in
           \sref{SEC:BC:NameMatch} and \tref{TBL:BC:NameMatch}.}
  \label{TBL:PW_BC_GFR_Match}
  \begin{tabular}{ @{} l r @{\qquad} l @{\qquad} c @{} }
    %%%%%%%%%%%%%%%%%%%%%%%%%%%%%%%%%%%%%%%%%%%%%%%%%%%%%%%%%%%%%%%%%%%%%%%%%%%%
    \toprule
    %%%%%%%%%%%%%%%%%%%%%%%%%%%%%%%%%%%%%%%%%%%%%%%%%%%%%%%%%%%%%%%%%%%%%%%%%%%%
    CGNS BC Option & GFR BC &
         \multirow[b]{2}{*}{Boundary Condition Description} &
         \multirow[b]{2}{*}{Section \#} \\ % Header top row
    (PW \texttt{"CAE Type"}) & Value & & \\ % Header bottom row
    %%%%%%%%%%%%%%%%%%%%%%%%%%%%%%%%%%%%%%%%%%%%%%%%%%%%%%%%%%%%%%%%%%%%%%%%%%%%
    \midrule
    %%%%%%%%%%%%%%%%%%%%%%%%%%%%%%%%%%%%%%%%%%%%%%%%%%%%%%%%%%%%%%%%%%%%%%%%%%%%
    Unspecified             & -9999 & Name match is attempted
                                    & (\ref{SEC:BC:NameMatch}) \\
    User Defined            & -9999 & Name match is attempted
                                    & (\ref{SEC:BC:NameMatch}) \\
    Axisymmetric Wedge      & -9999 & Name match is attempted
                                    & (\ref{SEC:BC:NameMatch}) \\
    Degenerate Line         & -9999 & Name match is attempted
                                    & (\ref{SEC:BC:NameMatch}) \\
    Degenerate Point        & -9999 & Name match is attempted
                                    & (\ref{SEC:BC:NameMatch}) \\
    Dirichlet               & -9999 & Name match is attempted
                                    & (\ref{SEC:BC:NameMatch}) \\
    Extrapolate             &    13 & Supersonic outflow
                                    & (\ref{SEC:BC:Outflow:Supersonic}) \\
    Farfield                &    21 & Characteristic inflow\slash outflow
                                    & (\ref{SEC:BC:Farfield:Characteristic}) \\
    General                 &    22 & General free-flow inflow\slash outflow
                                    & (\ref{SEC:BC:Farfield:General}) \\
    Inflow                  &     1 & General Inflow
                                    & (\ref{SEC:BC:Inflow:General}) \\
    Inflow Subsonic         &     2 & Subsonic Inflow
                                    & (\ref{SEC:BC:Inflow:Subsonic}) \\
    Inflow Supersonic       &     3 & Supersonic Inflow
                                    & (\ref{SEC:BC:Inflow:Supersonic}) \\
    Neumann                 & -9999 & Name match is attempted
                                    & (\ref{SEC:BC:NameMatch}) \\
    Outflow                 &    11 & General outflow
                                    & (\ref{SEC:BC:Outflow:General}) \\
    Outflow Subsonic        &    12 & Subsonic outflow
                                    & (\ref{SEC:BC:Outflow:Subsonic}) \\
    Outflow Supersonic      &    13 & Supersonic outflow
                                    & (\ref{SEC:BC:Outflow:Supersonic}) \\
    Symmetry Plane          &    80 & Symmetry boundary
                                    & (\ref{SEC:BC:Numerical:Symmetry}) \\
    Symmetry Polar          & -9999 & Name match is attempted
                                    & (\ref{SEC:BC:NameMatch}) \\
    Tunnel Inflow           & -9999 & Name match is attempted
                                    & (\ref{SEC:BC:NameMatch}) \\
    Tunnel Outflow          & -9999 & Name match is attempted
                                    & (\ref{SEC:BC:NameMatch}) \\
    Wall                    &    52 & Adiabatic wall
                                    & (\ref{SEC:BC:Wall:Adiabatic}) \\
    Wall Inviscid           &    50 & Inviscid\slash slip wall
                                    & (\ref{SEC:BC:Wall:Slip}) \\
    Wall Viscous            &    52 & Adiabatic Wall
                                    & (\ref{SEC:BC:Wall:Adiabatic}) \\
    Wall Viscous Heat Flux  &    52 & Adiabatic Wall
                                    & (\ref{SEC:BC:Wall:Adiabatic}) \\
    Wall Viscous Isothermal &    53 & Isothermal Wall
                                    & (\ref{SEC:BC:Wall:Isothermal}) \\
    \bottomrule
  \end{tabular}
\end{table}

%%%%%%%%%%%%%%%%%%%%%%%%%%%%%%%%%%%%%%%%%%%%%%%%%%%%%%%%%%%%%%%%%%%%%%%%%%%%%%%%

\subsubsubsection{Limitations}
\label{SEC:Files:Grid_Input:CGNS:Limitations}

Currently, there is no mechanism to read structured (\thatis $i,j,k$ ordered)
blocks\slash zones from CGNS grid files, so GFR only supports CGNS grid files
with unstructured zones.
%
Structured blocks that are written to a CGNS grid file as quad\slash hex-only
unstructured zones are fully supported.
%
When using Pointwise, converting structured blocks to quad\slash hex-only
unstructured zones is possible by choosing the option to combine all blocks when
using the regular (\thatis linear) CGNS export, which should result in an
unstructured grid file that can be used by GFR.
%
CGNS grids exported from Pointwise using the high-order option will always be
written as a single unstructured zone.
%
Already existing structured CGNS grid files must first be converted to
unstructured CGNS format using third-party software to use with GFR.

%%%%%%%%%%%%%%%%%%%%%%%%%%%%%%%%%%%%%%%%%%%%%%%%%%%%%%%%%%%%%%%%%%%%%%%%%%%%%%%%

GFR requires CGNS files to be HDF5 formatted.
%
ADF formatted CGNS files can be converted to HDF5 format using the
\texttt{cgnsconvert} utility that is included with the CGNS library.

%%%%%%%%%%%%%%%%%%%%%%%%%%%%%%%%%%%%%%%%%%%%%%%%%%%%%%%%%%%%%%%%%%%%%%%%%%%%%%%%

